\chapter{Аудит кандидатов IR-массового члена трансмутированной моды}
\label{ch:v10-discrete-resolution-transmuted-ir-mass-audit}

Предыдущий гейт показал, что условный полюс
\begin{equation}
 \Gamma_*^{(2)}(q)=q-e^{-64\pi^2/3}
\end{equation}
здоров, но не объяснил происхождение постоянного члена самоэнергии. Поэтому
проверены одиннадцать механизмов: bare mass, Higgs/Yukawa-конденсат,
асимптотически-свободный конденсат, KMS thermal mass, Lamb shift ванны,
кривизна, конечный объём, throughflow, portal splitting, наблюдаемый полюс
и формальная подстановка цели.

Каждый кандидат обязан одновременно иметь правильный тип $m^2$, внутренний
носитель, ненулевой IR-предел, выбранный родитель, точный показатель
$e^{-64\pi^2/3}$ и отсутствие target loading. Критериальная матрица имеет
размер $11\times6$, ранг $6$ и оценки
\begin{equation}
 (3,4,3,5,4,4,5,4,4,3,5).
\end{equation}
Полных проходов нет. KMS-масса и конечный объём внутренне допустимы, но их
коэффициенты не дают требуемой экспоненты. Наблюдаемый и формальный полюсы
кругово подставляют ответ. Остальные механизмы не имеют выбранного
коэффициента или нужного IR-предела.

Особый результат даёт проверка знака beta-функции. При условном
$b_{\rm AF}=-2$ и унаследованном $g^2=3/8$
\begin{equation}
 \log\frac{\mu}{k_{\rm IR}}
 =\frac{8\pi^2}{|b_{\rm AF}|g^2}
 =\frac{32\pi^2}{3},\qquad
 \frac{m_{\rm IR}^2}{\mu^2}=e^{-64\pi^2/3}.
\end{equation}
Это точный некруговой механизм обратной экспоненты. Но текущий носитель
имеет $b=+2$: знак и асимптотически-свободный сектор в нём не выведены.

\begin{theorem}[Единственный точный механизм отсутствует в носителе]
Ни один из одиннадцати текущих кандидатов не порождает требуемый IR-массовый
член по всем шести критериям. Условная замена $b=+2$ на $b=-2$ точно даёт
нужную иерархию, но является новой физической структурой.
\end{theorem}

\begin{nogocriterion}[Массовый член нельзя объявить ответом]
Bare counterterm, fitted pole или запись $m^2=e^{-64\pi^2/3}$ не закрывают
происхождение. Допустим только сектор, где отрицательный beta-коэффициент,
его носитель и связь конденсата с полюсом следуют из общего родителя.
\end{nogocriterion}

\begin{protocol}\begin{enumerate}
\item охват механизмов: \textbf{$11/11$};
\item полные проходы: \textbf{$0/11$};
\item точный условный AF-механизм: \textbf{$1/1$};
\item физическое происхождение: \textbf{$0/3$};
\item следующий гейт: \textbf{родитель отрицательного beta-знака}.
\end{enumerate}\end{protocol}

Машинный сертификат сохранён в
\path{s2t_v10_cell_birth_four_volume_nonequilibrium_bath_discrete_resolution_transmuted_mode_ir_mass_term_candidate_audit_gate_results.json}.